\newpage
\section{A Formal Definition of Data}
\label{sec:formal-definition}

The informal definitions reviewed in Section~\ref{sec:what-is-data} are
consistent with one another, but they are not formal: none specifies what
structure a collection of facts must have before it can be indexed, compared
for equality, or serialized. Mathematics offers a sharper handle on the same
idea. In mathematics, data are most naturally thought of as a
\emph{function}: a mapping from some index set to some value set.

\subsection{Primitive Notions}

\begin{definition}[Index Set and Value Set]
\label{def:index-value}
An \emph{index set} $I$ is a set whose elements serve as positions or
addresses. A \emph{value set} $V$ is a set whose elements serve as the
possible contents at a position.
\end{definition}

Neither $I$ nor $V$ is constrained further at this level of abstraction: $I$
may be finite or infinite, ordered or unordered, and $V$ may contain
elements of any kind. What matters is that both are fixed before a collection
of data is constructed.

\begin{definition}[Datum]
\label{def:datum}
A \emph{datum} is an element $v \in V$, drawn from some value set $V$,
designated as the content held at a single index.
\end{definition}

\subsection{The Function-Theoretic Definition}

\begin{definition}[Collection of Data]
\label{def:collection}
A \emph{collection of data} $D$ over index set $I$ and value set $V$ is a
total function
\[
  D : I \to V
\]
assigning to every index $i \in I$ exactly one value $D(i) \in V$.
\end{definition}

This definition requires three conditions to hold for $D$ to be well-defined.

\paragraph{Condition 1 (Totality).}
$D$ must be defined for every $i \in I$. A structure that leaves some indices
unassigned is a \emph{partial function}, and any operation that assumes totality
is unsound when applied to it without first checking the domain.

\paragraph{Condition 2 (Functionality).}
For every $i \in I$, $D(i)$ must denote a single, determinate value in $V$.
A mapping that assigns multiple candidate values to the same index---for
instance, two fields claiming the same byte offset under different schema
versions---is not a function and therefore not a single well-defined collection
of data.

\paragraph{Condition 3 (Codomain Consistency).}
Every value $D(i)$ must belong to the \emph{same} value set $V$ for all $i \in
I$, where $V$ is fixed in advance of constructing $D$.

\subsection{Equality}

The function-theoretic definition yields a precise criterion for equality that
differs importantly from mere representational identity.

\begin{definition}[Functional Equality]
\label{def:equality}
Two collections of data $D_1 : I_1 \to V_1$ and $D_2 : I_2 \to V_2$ are
\emph{equal}, written $D_1 = D_2$, if and only if $I_1 = I_2$, $V_1 = V_2$,
and $D_1(i) = D_2(i)$ for every $i \in I_1$.
\end{definition}

This is a stronger condition than byte-for-byte identity of an encoding. Two
structures can share an identical encoding while differing in $V$---as the
five-byte example in Section~\ref{sec:what-is-data} shows---and
Definition~\ref{def:equality} correctly classifies them as distinct collections
of data. Conversely, two collections may be equal under Definition~\ref{def:equality}
while differing in their encodings, because the encoding is not part of the
definition at all.

\subsection{Relation to Standard Mathematical Structures}

The function-theoretic definition deliberately generalizes several structures
often used informally as synonyms for ``data.''

\begin{description}
  \item[Sequence / Array.] The special case $I = \{0, 1, \dots, n-1\}$: a
        function from a finite, contiguous, totally ordered index set into $V$.
        This is the canonical definition of a sequence in the mathematical
        literature~\autocite{knuth1997}, and it is the formal starting point for
        arrays. Arrays are data over the specific index set that integer
        addressing supplies.
  \item[Tuple.] The special case where $I$ is small and fixed, and $V$ is
        permitted to vary per index ($D(i) \in V_i$). Tuples relax Condition~3
        deliberately and explicitly---a structurally different move from the
        \emph{implicit} relaxation that produces type-confusion defects
        (Section~\ref{sec:implications}).
  \item[Multiset.] Corresponds to the case where $I$ is considered only up to a
        permutation. A multiset is a quotient of the function-theoretic
        structure, not an alternative to it: the underlying function still
        exists; the multiset view simply ignores the ordering of $I$.
  \item[Relation / Table.] A relation with $k$ attributes over domains
        $V_1, \dots, V_k$ is, under this definition, a set of functions from
        $\{1,\dots,k\}$ into $V_1 \times \cdots \times V_k$. The relational
        model's primary key is a selection of indices into this function
        space.
\end{description}

The unifying observation is that each of these structures is a special case of
Definition~\ref{def:collection}, distinguished by the properties imposed on $I$,
$V$, or both. Moving between them is a matter of specifying or relaxing
constraints on those two sets.

\subsection{A Worked Example}

A weather station records temperature once per hour. Formally, this is a
function $D : \mathbb{N} \to \mathbb{R}$, where the index set is hour number
and the value set is temperature in degrees Celsius. The station's log for one
day is the restriction of $D$ to $I = \{0, 1, \dots, 23\}$, which is an
array. The highest temperature recorded is $\max_{i \in I} D(i)$, a computation
that is defined precisely because $I$ is finite and $D$ is total on it. If a
sensor fails and no reading is recorded at hour 14, the structure is no longer
total on $I$---it is a partial function---and any operation that assumes totality,
such as computing the daily average, is unsound unless the missing value is
explicitly handled.

This example illustrates why the three conditions matter in practice:
engineering decisions about how to handle missing values, what type a field
carries, and whether two records are the same all reduce to questions about
$I$, $V$, and the conditions in Definitions~\ref{def:collection}
and~\ref{def:equality}.